% !PRINTED 684
\weberpage{684}
% [source note] The same-parity restriction on x,y below applies only to odd m. For a fundamental discriminant -m the general condition is x congruent to m*y modulo 2; for m=4, x=0,y=1 represents i. The printed blanket restriction is retained.
Beide Determinanten haben also den Minimalwerth, und
beide Systeme sind Fundamentalsysteme.

Der Fall $\nu = 1$, den wir oben ausgeschlossen haben, tritt
ausser im Falle des rationalen Körpers, in dem nur die beiden
Einheiten $\pm 1$ existiren, im Falle des \gesperrt{imaginären quadratischen
 Körpers} ein. In diesem Falle haben wir im zwanzigsten
 Abschnitte die ganzen Zahlen des Körpers in der Form

\[
\frac{x + y\sqrt{-m}}{2}
\]

dargestellt, worin $-m$ eine Stammdiscriminante bedeutet, so
dass $m \equiv 0$ oder $\equiv 3\ (\mathrm{mod}\ 4)$ ist, und $x, y$ beide gerade oder
beide ungerade sind. Die Einheiten erhalten wir, wenn wir die
Gleichung

\[
x^2 + m y^2 = 4
\]

auf alle mögliche Arten in ganzen rationalen Zahlen lösen. Diese
Gleichung hat aber, wenn $m > 4$ ist, nur die zwei Lösungen
$x = \pm 2$, $y = 0$, und es giebt also in diesen Fällen, wie im
Körper der rationalen Zahlen, nur die zwei Einheiten $\pm 1$. Ist
$m = 4$, so findet man die vier Einheiten $\pm 1, \pm i$, und ist endlich
 $m = 3$, die sechs Einheiten

\[
\pm 1, \; \pm \frac{-1 + i\sqrt{3}}{2}, \; \pm \frac{-1 - i\sqrt{3}}{2}.
\]

In dem nächst einfachen Falle, $n = 2$, $\nu = 2$, d. h. im
reellen quadratischen Körper, fällt die Theorie der Einheiten
zusammen mit der im \S.\,127 des ersten Bandes behandelten
Theorie der Pell'schen Gleichung.

\begin{center}
\S.\,188.\\[4pt]
\gesperrt{Reducirte Zahlen.}
\end{center}

Ist $\alpha$ irgend eine von Null verschiedene ganze oder gebrochene
 Zahl des Körpers $\Omega$, so betrachten wir, wie bei den
Einheiten, das \gesperrt{System der conjugirten Logarithmen} von $\alpha$,
worunter wir, wie im \S.\,185, die \gesperrt{reellen Zahlen}

\weq{1}{\lambda_1 = \delta_1 \log|\alpha_1|, \; \lambda_2 = \delta_2 \log|\alpha_2|, \; \ldots, \; \lambda_\nu = \delta_\nu \log|\alpha_\nu|}

verstehen, die wir auch mit $\lambda_s(\alpha)$ bezeichnen. Ziehen wir das
Fundamentalsystem $\varepsilon_1, \varepsilon_2, \ldots, \varepsilon_{\nu-1}$ von Einheiten zu Rathe, so
%CONT>>
